% =====================================================
% 题型：余弦定理与正弦值选择题 (q_triangle_cosine_rule_sine.tex)
% =====================================================
% 难度：★★ (中档)
% =====================================================
\newcommand{\qTriangleCosineRuleSineStem}{在 \(\triangle ABC\) 中，角 \(A,B,C\) 的对边分别为 \(a,b,c\)。若 \(a=3,\ b=5\)，且 \(\cos C=\dfrac35\)，则 \(\sin A=\hfill（\quad）\)}

\newcommand{\qTriangleCosineRuleSineOptions}{%
  \mcoptions[4]{\(\dfrac35\)}{\(\dfrac45\)}{\(\dfrac{12}{25}\)}{\(\dfrac{24}{25}\)}%
}

\newcommand{\qTriangleCosineRuleSineAnswer}{C}

\newcommand{\qTriangleCosineRuleSineSolution}{%
  由 \(\cos C=\dfrac35\)，且 \(C\in(0,\pi)\)，得 \(\sin C=\dfrac45\)。\\
  根据余弦定理：
  \[
    c^2=a^2+b^2-2ab\cos C
    =9+25-2\cdot3\cdot5\cdot\frac35
    =16,
  \]
  所以 \(c=4\)。\\
  由正弦定理：
  \[
    \frac{a}{\sin A}=\frac{c}{\sin C},
  \]
  因而
  \[
    \sin A=\frac{a\sin C}{c}
    =\frac{3\cdot\frac45}{4}
    =\frac35.
  \]
  故选 A。%
}

\newcommand{\qTriangleCosineRuleSineDiagnosis}{%
  若已知余弦值后直接把 \(\sin C\) 写成 \(1-\cos C\)，说明基本恒等式混淆；另一高频断点是余弦定理求出 \(c^2\) 后忘记取正根。%
}

\newcommand{\qTriangleCosineRuleSineTeachingNotes}{%
  题目链条是“余弦值 → 正弦值 → 第三边 → 正弦定理转角”。适合作为正余弦定理联用的短链训练。%
}
